\chapter{Literature review}

\section{Behaviour trees}

Earlier behaviour tree research mainly focused on game artificial intelligence. Flórez-Puga et al. (2009) added queries and the reuse of existing behaviours to game behaviour trees, while Lim et al. (2010) used behaviour trees to create agents for the commercial real-time strategy game DEFCON. These studies show that behaviour trees were already used early on to organise game decisions containing multiple conditions and actions, and were not only suitable for simple demonstrations.

After that, related research began to further explain the execution method and theoretical features of behaviour trees. Marzinotto et al. (2014) proposed a unified behaviour tree framework for robot control, and Colledanchise and Ögren (2017, 2018) analysed behaviour trees from the perspectives of hybrid control and modularity. By the time Iovino et al. (2022) conducted their survey, behaviour tree research had already covered game artificial intelligence, robotics, automatic learning, planning and design tools. This development process shows that behaviour trees have gradually developed from a practical tool in games into a more general agent control method.

Before testing display optimisation functions, it is necessary to confirm that the behaviour tree itself can run correctly. Otherwise, it is impossible to judge whether problems in the test come from the display optimisation function or the behaviour tree execution system. In this project, Sequence can remember the current execution position, while Reactive Selector needs to recheck high-priority conditions. Random Selector should keep the same choice during one running period. Repeat and Wait must also clear their memory when restarted. The blackboard is used to share state between perception, conditions and actions, and Decorator is used to restrict or change branch execution without modifying the child node itself. The editor cannot treat node order and Decorator as purely visual information, because they affect execution results.

\section{Node-link trees, hierarchy and scale}

Node-link diagrams are suitable for displaying behaviour trees, because connections can directly represent parent-child relationships. The tidy tree layout proposed by Reingold and Tilford (1981) emphasises centred parent nodes, spacing within the same layer and separated subtrees, while Sugiyama et al. (1981) further discuss ordering within layers and node positions. These methods are very suitable for automatically generating a tidy tree, but they do not completely solve the problem of interactive editing, because users still want to drag nodes freely.

Whether a layout is better cannot be judged only by whether the picture is tidy. Purchase (1997) found that edge crossings had a large influence among the tested graph aesthetics. Ware et al. (2002) also showed that path continuity, crossings and branches affect path finding in node-link diagrams. The results of Purchase, Carrington and Allder (2002) further show that only some aesthetic metrics stably affect task performance, and that the semantic domain of the graph can also change the result.

After scale increases, nodes and connections become harder to observe at the same time. Ghoniem et al. (2005) compared node-link diagrams and matrix representations, and explained that results are affected by scale, density and task at the same time. SpaceTree added dynamic layout and selective display to node-link trees, also showing that representation methods and interaction methods need to be evaluated together (Plaisant et al., 2002). For behaviour trees, the node-link method is still the easiest way to express parent-child direction and sibling order.

Node-based interfaces are also used to organise other complex development processes. InstructPipe allows users to generate machine learning workflows and continue modifying them in a node graph (Zhou et al., 2025). Although its research object is not a behaviour tree, it also needs to handle node positions, connection relationships and complex workflow navigation, showing that these problems do not only exist in behaviour tree editors.

Layout design also needs to consider the structural requirements of different domains. Helmke et al. (2024) found that domain experts sometimes prioritise preserving technical structure in graphs instead of always reducing connection crossings or edge length. Behaviour trees also have their own structural requirements, such as parent-child direction and sibling execution order, so automatic layout cannot only pursue smaller occupied area. The systematic review of graph layout evaluation by Di Bartolomeo et al. (2024) also points out that the evaluation of layout algorithms needs to explain data, scale and the metrics used at the same time. This provides methodological support for this dissertation to use real behaviour trees, multiple node scales and several direct metrics.

\section{Local avoidance and position stability}

Interactive editing is different from one-time automatic layout. After a user drags a node, if the system rearranges the whole tree, overlap may disappear, but the original position memory will also be broken. Misue et al. (1995) summarise this as the problem of layout adjustment and the mental map. Dwyer et al. (2006) proposed that nodes need to be separated from each other while new positions should stay as close as possible to the original positions. Later research also explained that the original topology still needs to be preserved when constraints are added (Dwyer et al., 2009).

Behaviour tree cards have real width, height, text and labels, so they cannot be treated only as points without area. Gansner and Hu (2010) emphasise that overlap removal should preserve proximity relationships and the original graph shape as much as possible. This study provides direction for Smart Drag: the system should deal with new overlap, but should not take over the whole canvas because of it.

This project does not directly copy one complete algorithm. Instead, it limits the starting distance, refresh frequency, propagation layers and maximum affected range, and then uses parent-child and sibling relationships to decide which nodes move together. The system only preserves the parent-above-child relationship that was already true before dragging.

The experiment by Archambault and Purchase (2013) shows that whether preserving node positions is helpful depends on the specific task. Therefore, this project requires relayout to maintain relative relationships between nodes as much as possible and reduce unnecessary position changes, but it does not use this point alone to claim that the function improves editing experience.

\section{Semantic zoom and focus-context}

Semantic zoom and fisheye are both used to handle a large amount of information on a limited canvas, but they use different methods. Semantic zoom changes the displayed content according to zoom level, rather than only scaling all content proportionally. Pad++ showed that interface content can change with zoom level (Bederson et al., 1996), and Summers et al. (2003) also compared different semantic zoom methods in program visualisation. Hiding secondary fields when zoomed out and restoring details after zooming in is the main feature of this type of method.

Fisheye belongs to focus and context methods. It enlarges the area currently concerned by the user while keeping the surrounding structure. Furnas (1986) proposed using degree of interest to distinguish focus and surrounding content, and Sarkar and Brown (1994) applied a similar idea to graphical fisheye. Wang et al. (2019) further studied the structural distortion that fisheye magnification may cause in large graphs and proposed a structure-aware method.

These display methods are not necessarily suitable for all tasks. The review by Cockburn et al. (2009) points out that overview, zooming and focus-context methods each have their own advantages and costs. Büring et al. (2006) also did not find that fisheye must shorten task time in an experiment on small screens, although more participants preferred this display method. In addition, Jakobsen and Hornbæk (2013) explain that physical screen size, information space and display scale jointly affect interaction results.

These studies provide design support for semantic zoom, fisheye and screen size, but they cannot directly explain their effect in a large behaviour tree editor. Therefore, this dissertation treats semantic zoom and fisheye as two different functions and compares them separately under different tree sizes and screen sizes.

\section{Relationship emphasis and connection occlusion}

Large graphs do not necessarily need to hide nodes first. Ware and Bobrow (2005) used interactive emphasis to help users query node-link diagrams, while van Ham and Perer (2009) retained context from the target of interest. For behaviour trees, the advantage of this type of method is that it can keep the original spatial position while reducing the visual strength of unrelated branches. Therefore, Related Node Focus in this project uses borders and opacity to distinguish selected nodes, ancestors, descendants, sibling nodes and unrelated branches.

Connections may also be blocked by nodes or other edges. EdgeLens by Wong et al. (2003) bends connections near the focus to leave space for crowded areas, while hierarchical edge bundles by Holten (2006) organise a large number of connections along a hierarchy. The experiment by McGee and Dingliana (2012) explains that bundling may help observe overall clusters, but may reduce the speed and accuracy of tracing a single path. Reducing visual clutter therefore does not mean that every connection task becomes easier.

The studies above show that layout adjustment, semantic zoom, fisheye, highlighting and connection treatment can all reduce some display problems of large graphs. However, the graph types, tasks and evaluation conditions used in these studies are different, so they cannot directly explain the actual effect of these methods in a large behaviour tree editor. Behaviour trees also need to preserve parent-child structure, sibling execution order and node information, which makes them different from ordinary node graphs. Based on this research gap, this dissertation compares these display optimisation functions under different tree sizes and screen sizes, and checks their direct changes, suitable environments and main costs.
